% Section 8.1, from lectures/L08/cnn_work/README.md: the project the exercises in Chapters 8 to 10
% are written in.

\section{The CNN Pipeline}
\label{c8:sec:pipeline}
\code{cnn_work}, in \file{lectures/L08/cnn_work}, is the working copy of a C++17 convolutional
neural network (CNN), built up one layer at a time across Chapters~\ref{c8:ch:conv},
\ref{c9:ch:maxpool} and~\ref{c10:ch:flatten}. It's trained to classify four hand-drawn 4\,×\,4
pixel patterns: the digits 0 to 3.

The four training patterns, from \file{source/main.cpp}, where \code{1} is a filled pixel and
\code{0} is empty, are shown below:

\begin{console}
Digit 0:      Digit 1:      Digit 2:      Digit 3:
1 1 1 1       0 1 0 0       1 0 0 1       0 0 1 0
1 0 0 1       0 1 0 0       0 1 1 0       0 0 1 0
1 0 0 1       0 1 0 0       0 1 1 0       1 1 1 1
1 1 1 1       0 1 0 0       1 0 0 1       0 0 1 0
\end{console}

\subsection{Current status}
Not every layer is real yet: \file{source/ml/factory/factory.cpp} decides, per layer type, whether
to build the real class or a \code{Stub}, and each one not yet built is marked with a
\code{//! @todo} comment there:
\begin{itemize}
  \item \textbf{\code{Conv}} (\file{ml/conv_layer/conv.hpp} and \file{conv.cpp}): placeholder, see
    the exercises in \secref{c8:sec:exercises}.
  \item \textbf{\code{MaxPool}} (\file{ml/conv_layer/max_pool.hpp} and \file{max_pool.cpp}):
    placeholder, see the exercises in \secref{c9:sec:exercises}.
  \item \textbf{\code{Flatten}} (\file{ml/flatten_layer/flatten.hpp} and \file{flatten.cpp}):
    placeholder, see the exercises in \secref{c10:sec:exercises}.
  \item \textbf{\code{Dense}}: already real (built separately back in Chapter~\ref{c5:ch:dense}).
\end{itemize}
The project builds and runs as-is at every stage: with any layer still stubbed, \code{make} builds,
trains and prints predictions without error; they just won't be meaningful until every layer is
real, at the end of Chapter~\ref{c10:ch:flatten}. That's expected, not a bug.

\subsection{Structure}

\begin{console}
cnn_work/
├── Makefile
├── source/
│   ├── main.cpp
│   └── ml/
│       ├── cnn/cnn.cpp
│       ├── conv_layer/{conv,max_pool}.cpp
│       ├── dense_layer/dense.cpp
│       ├── flatten_layer/flatten.cpp
│       ├── factory/factory.cpp
│       ├── random/generator.cpp
│       └── utils.cpp
└── include/ml/
    ├── types.hpp, utils.hpp
    ├── act_func/{interface,relu,tanh,none,type}.hpp
    ├── cnn/{interface,cnn}.hpp
    ├── conv_layer/{interface,conv,max_pool,stub}.hpp
    ├── dense_layer/{interface,dense,stub}.hpp
    ├── flatten_layer/{interface,flatten,stub}.hpp
    ├── factory/{interface,factory,stub}.hpp
    └── random/{interface,generator}.hpp
\end{console}

\noindent Each layer type (\code{conv_layer}, \code{dense_layer}, \code{flatten_layer}) follows the
same pattern: an \code{Interface} that defines the contract, a concrete class that does the real
computation (\code{Conv}, \code{MaxPool}, \code{Dense}, \code{Flatten}), and a \code{Stub} variant
that only validates dimensions and returns \code{true} or \code{false} without computing anything.
\code{ml::factory::Factory} picks, layer by layer, which one to build; see \file{factory.cpp}.

\textbf{Data types} (\file{ml/types.hpp}): \code{Matrix1d = std::vector<double>},
\code{Matrix2d = std::vector<Matrix1d>}, \code{Matrix3d = std::vector<Matrix2d>}.

\subsection{\code{dense_layer::Dense}}
The one layer that's already real. A regular fully connected layer with weights and a bias.
\begin{itemize}
  \item \textbf{\code{feedforward()}:} \code{sum = bias[i] + Σ weight[i][j] * input[j]} per node,
    followed by the activation function.
  \item \textbf{\code{backpropagate()}:} computes
\begin{cppcode}
error[i] = (outputGradients[i] - myOutput[i])
         * activationFunction.delta(preActivationSum[i])
\end{cppcode}
    See the note in \secref{c8:sec:details} on how \code{outputGradients} is interpreted differently
    depending on whether the layer is the output layer or a hidden layer.
  \item \textbf{\code{optimize()}:} updates the weights and bias with
    \code{error * learningRate * input}.
\end{itemize}

\subsection{\code{ml::cnn::Cnn}: how the layers connect}
The constructor creates and connects the layers via the injected factory:
\begin{enumerate}
  \item A \code{Conv} layer and a \code{MaxPool} layer are added to the same list,
    \code{myConvLayers} (in that order).
  \item A \code{Flatten} layer is created based on the max pooling layer's output size.
  \item A \code{Dense} layer is created based on the flatten layer's output size.
\end{enumerate}

\textbf{\code{feedforward(input)}:} each layer is fed the previous layer's \code{output()}:

\begin{console}
input → Conv → MaxPool → Flatten → Dense → prediction
\end{console}

\textbf{\code{backpropagate(target)}:} goes the same path in reverse. The target vector
(\code{target}) is passed directly into the last dense layer's \code{backpropagate()}: that's where
\code{(outputGradients[i] - myOutput[i])} actually becomes \code{(target[i] - prediction[i])}. Each
preceding layer is then fed the next layer's \code{inputGradients()}: the dense layers in reverse
order, then the flatten layer, then \code{MaxPool}, and finally \code{Conv}.

\textbf{\code{optimize(learningRate)}:} updates every layer with the gradients already computed by
\code{backpropagate()}. The dense layers also need the previous layer's \code{output()} (the value
actually used during feedforward) as input for the weight update.

\textbf{\code{train(trainIn, trainOut, epochCount, learningRate)}:} pure stochastic gradient
descent (batch size 1, no mini-batching). The training order is reshuffled every epoch. If
\code{feedforward()}, \code{backpropagate()}, or \code{optimize()} fails for any example, training
is aborted immediately.

\textbf{\code{predict(input)}:} runs \code{feedforward()} and returns the last dense layer's
\code{output()}.

\subsection{Activation functions}

\begin{narrowtable}{@{}lll@{}}
  \thead{Function} & \thead{\code{output(x)}} & \thead{\code{delta(x)}} \\
  \midrule
  \code{Relu} & \code{x} if \code{x > 0}, otherwise \code{0} & \code{1} if \code{x > 0}, otherwise
    \code{0} \\
  \code{Tanh} & \code{tanh(x)}, range $(-1, 1)$ & \code{1 - tanh(x)²} \\
  \code{None} & \code{x} (identity) & \code{1} \\
\end{narrowtable}

\code{delta()} is always computed from the pre-activation sum (the value before the activation
function was applied), not from the already-activated output value.

\subsection{Random number generator}
\code{random::Generator} is a singleton that initializes \code{std::srand} once, on first use. It's
used both to randomize initial weight and bias values (via \code{randomStartVal()} in \code{Conv}
and \code{Dense}), and to reshuffle the training order every epoch.

\subsection{Building and running}

\begin{shell}
make
\end{shell}

\noindent This builds and runs the program directly (the \code{build} target followed by the
\code{run} target). \code{clean} removes the built binary.

\subsection{Implementation details worth knowing}
\label{c8:sec:details}
\textbf{\code{Dense::backpropagate()} reuses the same formula for the output layer and hidden
layers.} For the last dense layer, \code{outputGradients} is literally the target vector, so
\code{outputGradients[i] - myOutput[i]} becomes \code{target - prediction}, which is the deviation
the book defines. For earlier dense layers, what's passed in instead is the next layer's
already-computed \code{inputGradients()}, which is a gradient, not a target. Subtracting
\code{myOutput[i]} from it adds a term that isn't part of the gradient of any loss, so for a hidden
layer this is a deliberate simplification rather than the exact derivative. The networks in this
book still train, because the term is small next to the one that carries the signal, and keeping
one formula keeps the layer readable. Worth knowing before you compare this code against a
textbook, or against a framework, and find they don't agree.
